\begin{statement}[C.3.1; Proved: fourth order at the maximum]
On either root branch at $\tau=i$, let
$h=\mu-(1-\tau)/2$ and $\delta=\tau-i$. Then
\begin{align*}
 \tau&=i-\frac{3\pi i}{2}h^2-\frac{9\pi^2i}{8}h^4+O(h^6),\\
 h^2&=\frac{2i}{3\pi}\delta+\frac1{3\pi}\delta^2+O(\delta^3),\\
 \mu&=\frac{1-\tau}{2}
  \mathbin{\pm}\left(\frac{2i\delta}{3\pi}\right)^{1/2}
          \left(1-\frac{i}{4}\delta+O(\delta^2)\right).
\end{align*}
The square root is taken holomorphically on the chosen root chart.
Furthermore
\begin{align}\label{eq:cs-Li}
 \log N={}&\tfrac12\log2-\frac\pi2|h|^2
   +\frac{3\pi^2}{8}|h|^4
   -\frac{3\pi^2}{4}|h|^2\operatorname{Re}(h^2)
   -\frac{5G}{192}\operatorname{Re}(h^4)+O(|h|^6),\\
 D_{\rm part}={}&D_* -\frac2\pi|\delta|
  +\frac1\pi\bigl(|\delta|^2+|\delta|\operatorname{Im}\delta\bigr)
  +\frac{5G}{72\pi^3}\operatorname{Re}(\delta^2)
  +O(|\delta|^3).\nonumber
\end{align}
In particular the leading terms of K3 and Lead 9 are correct.
\end{statement}
\begin{proof}
Write $e(\tau)=\wp_\tau((1-\tau)/2)$ and
$E=\pi^2E_2/3$. The half-period Taylor series is
\[
 \wp_\tau((1-\tau)/2+h)
 =e+(3e^2-g_2/4)h^2+(3e^3-g_2e/4)h^4+O(h^6).
\]
The exact coefficients checked in K3 are
$e(i)=0$, $e'(i)=iG/(6\pi)$. This proves the claimed identity
$3e'(i)=(2i/\pi)G/4$. For clarity its closed form follows from
$e=\pi^2(\vartheta_2^4-\vartheta_4^4)/3$ and the heat equation:
if $A=\vartheta_2(i)^4=\vartheta_4(i)^4$, then
$G=4\pi^4A^2$ and $e'(i)=2\pi^3iA^2/3$.
The classical beta integral
$K(1/2)=\Gamma(1/4)^2/(4\sqrt\pi)$ and
$\vartheta_3(i)^2=2K(1/2)/\pi$ give
$G=16K(1/2)^4=\Gamma(1/4)^8/(16\pi^2)$.
The beta evaluation follows by changing variable in
$\int_0^\infty(1+x^4)^{-1/2}\,dx
=\tfrac14 B(1/4,1/4)$; it is the lemniscatic complete period.

Ramanujan's identities give $g_2'(i)=2iG$. Differentiating
$e(-1/\tau)=\tau^2 e(\tau)$ twice at $i$ gives
$e''(i)=3i e'(i)=-G/(2\pi)$. Inserting
$\delta=A_2h^2+A_4h^4+O(h^6)$ into the displayed half-period series
therefore gives
\[
 A_2=-3\pi i/2,\qquad
 A_4=\frac{g_2'A_2/4-e''A_2^2/2}{e'}=-9\pi^2i/8.
\]
The implicit-function theorem applies because $e'(i)\ne0$;
evenness in $h$ gives the stated remainder. Series inversion gives
the formulas for $h^2$ and $\mu$.

Let $H(h,\tau)=\log(\vartheta_3(h,\tau)/\eta(\tau))$. The exact
identity from the half-period translation is
$\log N=\operatorname{Re}H-\pi(\operatorname{Im}h)^2/
\operatorname{Im}\tau$. At $(0,i)$ its relevant Taylor data are
\[
 \begin{gathered}
 H_0=\tfrac12\log2,\quad H_\tau=0,\quad
 [h^2]H=-\pi/2,\quad [h^4]H=G/48,\\
 \partial_\tau[h^2]H=-\pi i/4-iG/(16\pi),\qquad
 H_{\tau\tau}=-G/(24\pi^2).
 \end{gathered}
\]
Indeed $\partial_h^2\log\vartheta_3=-\wp(h+(1-\tau)/2)-E$;
the heat equation at $h=0$ gives $H_\tau=ie/(4\pi)$, and
$E'(i)=\pi i/2-iG/(24\pi)$. These identities prove every entry.
After substitution $\tau=\tau(h)$, the pure $h^4$ coefficient of
$H$ is $-3\pi^2/8-5G/192$. Expanding
$\operatorname{Im}\tau=1-(3\pi/2)\operatorname{Re}(h^2)+O(|h|^4)$
in the Gaussian gives \eqref{eq:cs-Li}. Finally
\[
 |h|^2=\frac{2}{3\pi}|\delta|
          \left(1+\tfrac12\operatorname{Im}\delta+O(|\delta|^2)\right)
\]
and substitution into the quartic expression, followed by
multiplication by $6/\pi$, gives the formula for $D_{\rm part}$.
The script \path{Cspecial_expansion_checks.py} independently checks
these coefficients against theta evaluations; its output is retained.
\end{proof}
